\section{Discussion}

Temporal support in CMUR is scale-dependent under the tested source architecture. Medium-term direction is comparatively robust where the primary common two-sided classifications resolve, but restricting inference to verified source anchors reduces fitted-model precision. Most annual-change, turning and acceleration comparisons remain unresolved under both residual-amplitude scenarios and at every tested decision threshold. Together, these findings distinguish support for sustained direction from support for local temporal features of the same annually disseminated indicator. An unresolved comparison identifies an unmet support criterion, rather than an absence of real temporal change.

The medium-term result needs to be read jointly with its precision consequence. The primary comparison yields 12 agreements, no resolved-class disagreements and 15 unresolved Member-State cases. Unresolved directions increase from eight in the released annual analysis A to 15 in the source-anchor analysis B1. Classical 95\% OLS intervals widen in all 27 Member States, with a median paired B1/A width ratio of 1.858 (Table~\ref{tab:mediumterm}; Supplementary Tables~\ref{tab:S1}--\ref{tab:S2}). Thus, source-anchor restriction creates no resolved directional reversal under this rule, while leaving fewer directions supported and materially reducing model-based precision. Agreement for EU27\_2020 is consistent with that directional stability. The secondary common directional rule, yielding 15 agreements, no resolved disagreements and 12 unresolved Member-State comparisons, corroborates the directional finding under a different common support criterion. Neither comparison establishes equal slopes, and interval widths describe precision conditional on the classical linear OLS model, rather than the full uncertainty of the indicator.

The short-window findings persist when the residual amplitude is reduced. At $\tau=0.95$, the secondary Brownian-motivated $1/\sqrt{2}$ amplitude sensitivity B resolves 36 additional Member-State comparisons relative to the primary unscaled four-year-residual stress scenario U: 19 annual-change, 11 turning and six acceleration instances, with none losing resolution. Nevertheless, unresolved comparisons remain the large majority for each metric. This persistence extends to all seven tested thresholds from 0.80 to 0.99 under both scenarios (Supplementary Table~\ref{tab:S3}), so majority unresolvedness does not hinge on selecting exactly 0.95. The scenarios test sensitivity to residual amplitude; neither identifies the unobserved true two-year, odd-year error distribution or provides a probabilistic bound on it. Additional resolutions under B therefore indicate sensitivity of support to the assumed amplitude, rather than greater accuracy.

Source-aware analysis can also produce a resolved reversal at more permissive thresholds. For BE in 2019, annual change under B is resolved as positive at 0.80 and 0.85, whereas the released-series change is negative. The comparison becomes unresolved at 0.90 and the higher tested thresholds. This adverse case qualifies the reference-threshold absence of resolved disagreement: source-aware sensitivity can alter supported direction as well as remove resolution. The broad unresolvedness result and the individual reversal describe different aspects of the evidence and both matter for interpretation.

Continuous weakest-variant support adds gradation to the threshold classifications (Supplementary Table~\ref{tab:S4}). Scenario-concordant medians are higher under B for all three metrics, with a comparatively small difference for acceleration. These descriptive conditional summaries show variation in support that a single binary threshold compresses. Because the concordant subsets can differ between scenarios, their medians are not a paired-subset comparison. Read alongside the complete threshold grid, they help distinguish graded evidence from a classification artifact at one cutoff, without assigning inferential probabilities to the claims.

The midpoint diagnostics narrow an alternative explanation for acceleration unresolvedness. For anchors $a,b$ and midpoint $m=(a+b)/2$, the midpoint-centred second difference $b-2m+a$ is exactly zero. This geometry matters to the deterministic vote, but removing that vote resolves zero additional acceleration claims in the unscaled reference analysis, both when retaining the literal minimum-three requirement and under the separately labelled exploratory minimum-two stochastic-only convention (Supplementary Table~\ref{tab:S5}A). Midpoint geometry alone is therefore insufficient to explain the boundary. The result concerns resolution under these diagnostic conventions; it does not partition unresolvedness into causal contributions.

The multi-span evidence likewise limits what can be inferred about residual amplitude. In 76 same-geography/same-centre pairs across 26 countries, median absolute residuals are 0.525 percentage points for four-year spans and 0.750 for eight-year spans; the eight-year value is larger in 54 cases, the four-year value in 20, and two are equal (Supplementary Table~\ref{tab:S5}B). Larger residuals are more frequent at the longer span in this sample, but the dependent descriptive windows establish neither casewise monotonicity nor a general span-scaling relationship. They do not validate the Brownian amplitude assumption or identify a two-year calibration. The amplitude scenarios consequently remain explicit sensitivity conditions.

The closest same-object temporal comparator is \textcite{Haider2026}, whose EU-level reporting-year sensitivity addresses trend estimation, nominal precision and forecast interpretation while recognizing CMUR's biennial source. Its aggregate evidence and the present EU directional agreement overlap substantively, although their periods, samples and estimands differ. The additional empirical contribution here is the Member-State-wide claim-support map: common-rule medium-term direction and precision linked to annual-change, turning and acceleration support, paired amplitude sensitivities, threshold robustness and diagnostics of the boundary. This extends the scale and claim classes of the comparison rather than establishing the source-frequency premise anew. \textcite{Christis2023} and \textcite{Nuss2024Nowcasting} also explicitly recognize the lower-frequency waste input and estimated intervening years in their trend, target and nowcasting work.

This contribution also builds on established temporal-disaggregation and mixed-frequency practice \parencite{ESS2018,Quilis2018,Mattera2023}. Constructing, benchmarking or reconstructing a higher-frequency series and assessing support for temporal interpretations of an already released indicator are related but distinct tasks. The present evidence addresses the latter for CMUR. Environmental observation-frequency research supplies an adjacent empirical rationale for examining temporal inference \parencite{CortesGuzman2026}; its quantitative results do not determine the CMUR support boundary.

The scope of that boundary is the 27 Member States and EU27\_2020 over the fixed 2010--2022 medium-term and through-2022 short-window designs. Current public source-component vintages are asynchronous relative to released CMUR, motivating use of released values at verified anchors rather than synthetic reconstruction from mismatched components. The unobserved odd-year errors leave U and B as tested sensitivities, not identified uncertainty. Consensus requirements and support thresholds are explicit design conditions, with robustness assessed over the stated grid; transfer to other indicators or future CMUR vintages requires empirical replication. Within this scope, the evidence supports distinguishing published cadence from support for the temporal claim being made. Medium-term directional interpretation is better supported than local annual-change, turning and acceleration interpretation under the tested design, providing a concrete basis for more discriminating use of CMUR in temporal monitoring.
